In all our discussion so far we have implicitly assumed that the target, which often is the proton, see \cref{sec:diverse}, is a massless particle.
However, we can relax that condition and consider a target with finite mass $M$, which yields \enquote{Target Mass Corrections} (TMC).
We refer to Refs.~\cite{Schienbein:2007gr,Ruiz:2023ozv} for a detailed review on the topic and instead focus here on the practical consequences.

Our master formula, \cref{eq:fact1}, gives the so-called leading twist expression, which is dominant with respect to higher twist contributions, which are summarized as $\order{\LQCD^2/Q^2}$.
While these higher twist terms are in principle not accessible in pQCD as PDFs are by definition leading order objects, they have recently gained renewed interest by PDF fitting groups~\cite{Cerutti:2025yji,Ball:2025xtj,Harland-Lang:2025wvm}.
TMCs are formally a part of these higher twist contributions, but since they can be computed systematically and consistently they are typically included in modern PDF extractions.

In practice, TMC amount to a shift in kinematics and a reweighting.
We define
\begin{equation}
    r = \sqrt{1 + \frac{4x^2 M^2}{Q^2}} \label{eq:TMCr}
\end{equation}
with which we define the so-called Nachtmann variable~\cite{Nachtmann:1973mr}
\begin{equation}
    \xi_N = \frac{2x}{1 + r}\,. \label{eq:Nachtmann}
\end{equation}
Eventually, we find for the structure function $F_2$
\begin{equation}
    F_2^\text{TMC}(x,Q^2) = \frac{x^2}{\xi_N^2 r^3} F_2^\text{noTMC}(\xi_N, Q^2) + \frac{6M^2 x^3}{Q^2 r^4}h_2(\xi_N,Q^2) + \frac{12M^4 x^4}{Q^4 r^5}g_2(\xi_N,Q^2) \label{eq:TMC}
\end{equation}
where $F_2^\text{noTMC}$ refers to the setup we have discussed so far and
\begin{align}
    h_2(\xi_N,Q^2) &= \left(\tilde h_2 \otimes F_2^\text{noTMC}(Q^2)\right)(\xi_N) &\text{with}\quad \tilde h_2(z) &= \frac z {\xi_N}\,, \label{eq:TMCh2}\\
    g_2(\xi_N,Q^2) &= \left(\tilde g_2 \otimes F_2^\text{noTMC}(Q^2)\right)(\xi_N) &\text{with}\quad \tilde g_2(z) &= 1 - z\,. \label{eq:TMCg2}
\end{align}
For the other structure functions $F_1$, $F_L$, and $F_3$ one can find similar equations~\cite{Schienbein:2007gr}.
TMC are relevant for large $x$ and/or small $Q^2$ as follows from \cref{eq:TMCr,eq:Nachtmann}.
Examining \cref{eq:TMC} we note that the structure function at the experimentally measured Bjorken-$x$ corresponds to first approximation to a structure function computed at the Nachtmann variable $\xi_N$.
This is a simple consequence from the reduced phase space, which is available for massive particles.

Note that \cref{eq:TMCh2,eq:TMCg2} are convolutions, i.e.\ they require the evaluation of $F_2^\text{noTMC}$ at intermediate momentum fraction, which can be challenging in practice.
Thus one often applies the following approximation~\cite{Schienbein:2007gr}
\begin{equation}
    F_2^\text{TMC}(x,Q^2) \approx \frac{x^2}{\xi_N^2 r^3} F_2^\text{noTMC}(\xi_N, Q^2) \left(1 + \frac{6 M^2 x \xi_N}{Q^2 r}(1-\xi_N)^2\right)
\end{equation}
which only requires a single evaluation.

TMC introduce an explicit dependence on the target mass $M$ into \cref{eq:fact1} and capture a consistent part of the higher twist contributions $\order{\LQCD^2/Q^2}$.
TMCs are independent from the requested perturbative accuracy or the exchanged boson since they refer to a kinematic correction and thus apply to every calculation.
While they don't interfere directly with our previous content they must be applied separately on top of everything as a final operation after all other corrections have been taken into account.
Note that \cref{box:DISreq} should read more precisely
\begin{conclusionbox}[label={box:DISreqM}]{Improved DIS requirements}
    We require $Q^2 \gg M^2$ (deep) and $W^2 \gg M^2$ (inelastic) to call a lepton-hadron scattering DIS.
\end{conclusionbox}
